% GENERATED FILE. Edit canonical lessons, then run npm run generate:books.
% canonical-source-hash: fnv1a64:f283e069608fea13
% canonical-lessons: ML-C330-muti, ML-C330-bharani, ML-C330-kutam, ML-C330-cakku, ML-C330-svarnnam

\chapter{Lid, Jar, Water pot, Sack, Gold}
\label{ch:ml-lid-jar-water-pot-sack-gold}
\hlchaptermodality{330}

\begin{chapteropening}
\textbf{By the end of this chapter:} \emph{I can say a lid, a jar, a water pot, a sack and gold.}

Complete the last lesson of chapter 330: I can say a lid, a jar, a water pot, a sack and gold.
\end{chapteropening}

\section[\texorpdfstring{\ml{മൂടി} (mūṭi)}{മൂടി (mūṭi)}]{\ml{മൂടി} (mūṭi) — a lid}

\label{lesson:ML-C330-muti}



\begin{quote}
Before the new one: say the Malayalam for football, then the Malayalam for cricket.
\end{quote}

\subsection*{You'll want to know: \ml{മൂടി}}
\textbf{\ml{മൂടി}} — \emph{mūṭi} — \textquotedblleft{}a lid\textquotedblright{}.

\begin{cousinweb}
\textbf{\ml{മൂടി}} (\emph{mūṭi}) is native Dravidian, the noun of the verb \textquotedblleft{}to cover\textquotedblright{}; DEDR 5034 lists the root with Tamil \emph{mūṭu}.
\end{cousinweb}

\begin{grammarlens}[title={What you've built}]
One more word for this chapter.
\end{grammarlens}

\subsection*{Your turn}
\begin{itemize}
\raggedright
  \item \emph{Say it:} \emph{mūṭi}
  \item \emph{Say it:} \emph{mūṭi}, once more
  \item \emph{Say it:} say \emph{mūṭi}
  \item \emph{Recall:} say the Malayalam for football, then the Malayalam for cricket, then say \emph{mūṭi} again
\end{itemize}

\begin{tcolorbox}[breakable,colback=teal!4,colframe=teal!35!black,title={Before you move on}]
What does \ml{മൂടി} mean? (\textquotedblleft{}A lid\textquotedblright{}.) Say it once more.
\end{tcolorbox}
\section[\texorpdfstring{\ml{ഭരണി} (bharaṇi)}{ഭരണി (bharaṇi)}]{\ml{ഭരണി} (bharaṇi) — a jar}

\label{lesson:ML-C330-bharani}



\begin{quote}
Before the new one: say the Malayalam for cricket, then the Malayalam for a lid.
\end{quote}

\subsection*{You'll want to know: \ml{ഭരണി}}
\textbf{\ml{ഭരണി}} — \emph{bharaṇi} — \textquotedblleft{}a jar\textquotedblright{}.

\begin{cousinweb}
\textbf{\ml{ഭരണി}} (\emph{bharaṇi}) has the shape of Sanskrit \emph{bharaṇī}, \textquotedblleft{}bearing\textquotedblright{}, also the name of a lunar asterism (CDIAL 9396); whether the jar's name comes from it is uncertain.
\end{cousinweb}

\begin{grammarlens}[title={What you've built}]
One more word for this chapter.
\end{grammarlens}

\subsection*{Your turn}
\begin{itemize}
\raggedright
  \item \emph{Say it:} \emph{bharaṇi}
  \item \emph{Say it:} \emph{bharaṇi}, once more
  \item \emph{Say it:} say \emph{bharaṇi}
  \item \emph{Recall:} say the Malayalam for cricket, then the Malayalam for a lid, then say \emph{bharaṇi} again
\end{itemize}

\begin{tcolorbox}[breakable,colback=teal!4,colframe=teal!35!black,title={Before you move on}]
What does \ml{ഭരണി} mean? (\textquotedblleft{}A jar\textquotedblright{}.) Say it once more.
\end{tcolorbox}
\section[\texorpdfstring{\ml{കുടം} (kuṭaṁ)}{കുടം (kuṭaṁ)}]{\ml{കുടം} (kuṭaṁ) — a water pot}

\label{lesson:ML-C330-kutam}



\begin{quote}
Before the new one: say the Malayalam for a lid, then the Malayalam for a jar.
\end{quote}

\subsection*{You'll want to know: \ml{കുടം}}
\textbf{\ml{കുടം}} — \emph{kuṭaṁ} — \textquotedblleft{}a water pot\textquotedblright{}.

\begin{cousinweb}
\textbf{\ml{കുടം}} (\emph{kuṭaṁ}) is native Dravidian, one word with Tamil \emph{kuṭam}, Kannada \emph{koḍa} and Telugu \emph{kuḍaka} (DEDR 1651). DEDR 1651 compares Sanskrit \emph{kūṭa}.
\end{cousinweb}

\begin{grammarlens}[title={What you've built}]
One more word for this chapter.
\end{grammarlens}

\subsection*{Your turn}
\begin{itemize}
\raggedright
  \item \emph{Say it:} \emph{kuṭaṁ}
  \item \emph{Say it:} \emph{kuṭaṁ}, once more
  \item \emph{Say it:} say \emph{kuṭaṁ}
  \item \emph{Recall:} say the Malayalam for a lid, then the Malayalam for a jar, then say \emph{kuṭaṁ} again
\end{itemize}

\begin{tcolorbox}[breakable,colback=teal!4,colframe=teal!35!black,title={Before you move on}]
What does \ml{കുടം} mean? (\textquotedblleft{}A water pot\textquotedblright{}.) Say it once more.
\end{tcolorbox}
\section[\texorpdfstring{\ml{ചാക്ക്} (cākkŭ)}{ചാക്ക് (cākkŭ)}]{\ml{ചാക്ക്} (cākkŭ) — a sack}

\label{lesson:ML-C330-cakku}



\begin{quote}
Before the new one: say the Malayalam for a jar, then the Malayalam for a water pot.
\end{quote}

\subsection*{You'll want to know: \ml{ചാക്ക്}}
\textbf{\ml{ചാക്ക്}} — \emph{cākkŭ} — \textquotedblleft{}a sack\textquotedblright{}.

\begin{cousinweb}
\textbf{\ml{ചാക്ക്}} (\emph{cākkŭ}) $\leftarrow$ English \emph{sack}, a word that goes back through Latin and Greek to a Semitic word for sackcloth.
\end{cousinweb}

\begin{grammarlens}[title={What you've built}]
One more word for this chapter.
\end{grammarlens}

\subsection*{Your turn}
\begin{itemize}
\raggedright
  \item \emph{Say it:} \emph{cākkŭ}
  \item \emph{Say it:} \emph{cākkŭ}, once more
  \item \emph{Say it:} say \emph{cākkŭ}
  \item \emph{Recall:} say the Malayalam for a jar, then the Malayalam for a water pot, then say \emph{cākkŭ} again
\end{itemize}

\begin{tcolorbox}[breakable,colback=teal!4,colframe=teal!35!black,title={Before you move on}]
What does \ml{ചാക്ക്} mean? (\textquotedblleft{}A sack\textquotedblright{}.) Say it once more.
\end{tcolorbox}
\section[\texorpdfstring{\ml{സ്വർണ്ണം} (svarṇṇaṁ)}{സ്വർണ്ണം (svarṇṇaṁ)}]{\ml{സ്വർണ്ണം} (svarṇṇaṁ) — gold}

\label{lesson:ML-C330-svarnnam}



\begin{quote}
Before the new one: say the Malayalam for a water pot, then the Malayalam for a sack.
\end{quote}

\subsection*{You'll want to know: \ml{സ്വർണ്ണം}}
\textbf{\ml{സ്വർണ്ണം}} — \emph{svarṇṇaṁ} — \textquotedblleft{}gold\textquotedblright{}.

\begin{cousinweb}
\textbf{\ml{സ്വർണ്ണം}} (\emph{svarṇṇaṁ}) $\leftarrow$ Sanskrit \emph{svarṇa}, \textquotedblleft{}gold\textquotedblright{}, taken over whole as a tatsama, with Malayalam's noun ending \emph{-ṁ}.
\end{cousinweb}

\begin{grammarlens}[title={What you've built}]
One more word for this chapter.
\end{grammarlens}

\subsection*{Your turn}
\begin{itemize}
\raggedright
  \item \emph{Say it:} \emph{svarṇṇaṁ}
  \item \emph{Say it:} \emph{svarṇṇaṁ}, once more
  \item \emph{Say it:} say \emph{svarṇṇaṁ}
  \item \emph{Recall:} say the Malayalam for a water pot, then the Malayalam for a sack, then say \emph{svarṇṇaṁ} again
\end{itemize}

\begin{tcolorbox}[breakable,colback=teal!4,colframe=teal!35!black,title={Before you move on}]
What does \ml{സ്വർണ്ണം} mean? (\textquotedblleft{}Gold\textquotedblright{}.) Say it once more.
\end{tcolorbox}
